\subsection{分解定理}

定理1（分解定理）. --- 设 \((x_i)_{1 \leq i \leq p}\) 与 \((y_j)_{1 \leq j \leq q}\) 是格序群 \(G\) 中正元素的两个有限序列，满足 \(\sum_{i=1}^{p} x_i = \sum_{j=1}^{q} y_j\) ；则存在正元素的一个双重序列 \((z_{ij})_{1 \leq i \leq p, 1 \leq j \leq q}\) （属于 \(G\) ），使得 \(x_i = \sum_{j=1}^{q} z_{ij}\) （对所有 \(i\) ），并且 \(y_j = \sum_{i=1}^{p} z_{ij}\) （对所有 \(j\) ）。

\begin{enumerate}
\def\labelenumi{\arabic{enumi})}
\tightlist
\item
  我们首先对 \(p = q = 2\) 的情形证明该定理。设 \(x, x', y, y'\) 是 \(G\) 的正元素，满足 \(x + x' = y + y'\) ，并令 \(a = \sup (0, x - y')\) 。由于
\end{enumerate}

\[
x - y ^ {\prime} = y - x ^ {\prime}
\]

小于 \(x\) 且小于 \(y\) ，由此可得 \(b = x - a\) 与 \(c = y - a\) 是正的， \(d = a - (x - y')\) 也是正的。此外我们还有

\[
x = a + b, x ^ {\prime} = c + d, y = a + c \quad \text { and } \quad y ^ {\prime} = b + d.
\]

\begin{enumerate}
\def\labelenumi{\arabic{enumi})}
\setcounter{enumi}{1}
\tightlist
\item
  现在我们证明：若定理对 \(p < m\) 与 \(q = n\) ( \(m > 2\) , \(n \geqslant 2\) )成立，则它对 \(p = m\) 与 \(q = n\) 也成立。根据假设，我们有 \(x_{m} + \sum_{i=1}^{m-1} x_{i} = \sum_{j=1}^{n} y_{j}\) 。由于定理对 \(p = 2\) 与 \(q = n\) 成立，故存在两个有限序列 \((z_{j}')\) , \((z_{j}''')\) （各由 \(n\) 个正项组成），使得 \(\sum_{i=1}^{m-1} x_{i} = \sum_{j=1}^{n} z_{j}'\) , \(x_{m} = \sum_{j=1}^{n} z_{j}''\) ，并且 \(y_{j} = z_{j}' + z_{j}''\) （对 \(1 \leqslant j \leqslant n\) ）。另一方面，由于定理对 \(p = m-1\) 与 \(q = n\) 成立，存在一个双重序列 \((u_{ij})_{1 \leqslant i \leqslant m-1, 1 \leqslant j \leqslant n}\) ，使得 \(x_{i} = \sum_{j=1}^{n} u_{ij}\) （对 \(1 \leqslant i \leqslant m-1\) ），并且 \(z_{j}' = \sum_{i=1}^{m-1} u_{ij}\) （对 \(1 \leqslant j \leqslant n\) ）。令
\end{enumerate}

\[
z _ {i j} = u _ {i j} \quad \text { for } \quad 1 \leqslant i \leqslant m - 1, \quad \text { and } \quad z _ {m j} = z _ {j} ^ {\prime \prime} \quad (1 \leqslant j \leqslant n),
\]

我们便确实得到一个满足定理条件的双重序列。3) 通过交换 \(x_{i}\) 与 \(y_{j}\) ，我们以同样的方式看到，若定理对 \(p = m\) 与 \(q < n\) ( \(m \geqslant 2\) , \(n > 2\) )成立，则它对 \(p = m\) 与 \(q = n\) 也成立。于是定理通过从 \(p = q = 2\) 情形出发的双重归纳法得到证明，因为只要 \(p \leqslant 1\) 或 \(q \leqslant 1\) ，它是平凡成立的。

推论。 --- 设 \(y, x_1, x_2, \ldots, x_n\) 是 \(n + 1\) 个正元素（属于 \(G\) ），满足 \(y \leqslant \sum_{i=1}^{n} x_i\) ；则存在 \(n\) 个正元素 \(y_i (1 \leqslant i \leqslant n)\) ，使得 \(y_i \leqslant x_i\) 且 \(y = \sum_{i=1}^{n} y_i\) 。

只需对序列 \((x_{i})\) 以及由两个元素 \(y\) 与 \(z = \left(\sum_{i=1}^{n} x_{i}\right) - y\) 组成的序列应用定理1。

